\chapter{Conclusiones}
\label{sec:conclusiones}

Tras el diseño, simulación y análisis de los moduladores y demoduladores de AM y BLU, se presentan las siguientes
conclusiones:

\section{Módem AM con Frecuencia Intermedia}
\begin{itemize}
    \item La simulación en Multisim permitió validar de manera práctica los conceptos teóricos de modulación con
          portadora. Se comprobó que el espectro posee las tres componentes esperadas (portadora y bandas laterales).
    \item En la conversión a FI, se demostró que el método superheterodino ($f_{\text{LO}} > f_c$) es preferible frente al
          subheterodino ($f_{\text{LO}} < f_c$) en términos de diseño de filtros. En el superheterodino, las componentes
          de suma del mezclador se sitúan en $250\kHz$, alejándose de la banda de FI de $50\kHz$. En cambio, en el
          subheterodino, estas espurias aparecen en $150\kHz$, lo que exigiría un filtro con una transición de banda de
          atenuación mucho más abrupta.
    \item Respecto a la demodulación, se contrastaron las dos tecnologías:
        \begin{itemize}
            \item El detector de envolvente es una opción simple y económica. Sin embargo, su desempeño está acotado por la
                  constante de tiempo del circuito RC, y es ineficiente ante esquemas con portadora suprimida o sobremodulados.
            \item El detector sincrónico ofrece una demodulación lineal y limpia de distorsiones, incluso para modulación
                  con portadora suprimida. Su desventaja radica en la complejidad circuital que implica recuperar la portadora
                  en sincronía de fase y frecuencia.
        \end{itemize}
\end{itemize}

\section{Módem BLU por Método de Desfase}
\begin{itemize}
    \item Este método es una alternativa elegante para generar BLU sin recurrir a filtros de banda ultra-selectivos
          de cristal. Esto se logra mediante la cancelación por desfase en cuadratura de la banda lateral no deseada.
    \item A través de la simulación y el análisis matemático, se corroboró la alta sensibilidad del modulador ante
          las imperfecciones de fase. Un desfasaje de $80^\circ$ (desviación de $10^\circ$ con respecto al valor ideal de
          $90^\circ$) provoca que la banda lateral no deseada se atenúe únicamente $-21.16\text{ dB}$ con respecto a la deseada.
          Este rechazo es insuficiente para la mayoría de los estándares modernos de comunicaciones.
    \item Variaciones en la amplitud de la banda base (Punto A) demostraron que el modulador mantiene su linealidad,
          afectando únicamente la potencia de la señal de salida sin alterar la supresión de la banda lateral.
    \item Para lograr altos niveles de rechazo de banda lateral en la práctica (típicamente mayores a $-40\text{ dB}$),
          los componentes analógicos de desfase deben estar muy bien apareados o bien el proceso de desfase debe
          ser implementado de forma digital mediante un transformador de Hilbert.
\end{itemize}
